\section{Experiments and Results}
\label{sec:results}

%% ============================================================================
%% DATA PROVENANCE (for authors; remove before submission)
%% Canonical model:  EfficientNetB0 + SMOTE, served preprocessed .npy arrays
%%                    (notebook: temp_paper_stuff/SMOTE.ipynb).
%% Test set:          366 fundus images (held out, never resampled).
%% Model selection:   best validation quadratic weighted kappa (QWK).
%% Operating point:   class label = argmax of the 5-way softmax output.
%%
%% PLACEHOLDERS marked "[to supply]" must be filled once the
%% SMOTE + Focal-loss experiment is run to completion (test metrics,
%% confusion matrix, AUC, Grad-CAM). They live only in
%% Sec.~\ref{subsec:imbalance_results}.
%%
%% Figures are intentionally left as commented \includegraphics stubs;
%% uncomment and supply the exported PNG/PDF for the 366-image run.
%% ============================================================================

The trained diabetic-retinopathy (DR) grading model is evaluated on a held-out
test set of $366$ fundus images that was excluded from all training and from the
synthetic oversampling procedure. Unless stated otherwise, the reported model is
the EfficientNetB0 classifier trained with SMOTE-balanced training data, and the
predicted grade for each image is obtained as the $\arg\max$ of the five-way
softmax output. The model checkpoint used for evaluation is the one that
maximised the validation quadratic weighted kappa (QWK). Performance is
characterised along three complementary axes: overall and per-class
classification quality (Section~\ref{subsec:classification_results}), the spatial
plausibility of the model's evidence via Grad-CAM
(Section~\ref{subsec:xai_results}), and the effect of the class-imbalance
strategy on both of these (Section~\ref{subsec:imbalance_results}).

\subsection{Experimental Setup and Evaluation Protocol}
\label{subsec:eval_protocol}

The classifier is an ImageNet-pretrained EfficientNetB0 backbone followed by a
classification head (global average pooling $\rightarrow$
\texttt{Dense(256)} $\rightarrow$ batch normalisation $\rightarrow$
\texttt{Dropout(0.5)} $\rightarrow$ \texttt{Dense(128)} $\rightarrow$
\texttt{Dropout(0.4)} $\rightarrow$ \texttt{Dense(5)}), totalling approximately
$4.41$M parameters. Training was conducted in two phases: a frozen-backbone phase
in which only the head was optimised, followed by a fine-tuning phase in which the
upper backbone layers (from layer $100$ onward, batch-normalisation layers kept
frozen) were unfrozen. Both phases used the Adam optimiser with a cosine-decay
learning-rate schedule, and both monitored the validation QWK for checkpointing
and early stopping.

Class imbalance was addressed by applying SMOTE to the \emph{training split only},
oversampling every class up to the size of the majority training class so that the
five grades were represented equally during training. The validation and test
splits retain the natural, imbalanced grade distribution and were never resampled.
The resulting class distributions are summarised in
Table~\ref{tab:data_distribution}.

\begin{table}[H]
\centering
\caption{Class distribution across the data splits. SMOTE was applied to the
training split only; the validation and test splits retain the original
(imbalanced) grade frequencies.}
\label{tab:data_distribution}
\begin{tabular}{lccc}
\toprule
Class & Train (post-SMOTE) & Validation & Test \\
\midrule
No DR            & 1{,}434 & 172 & 199 \\
Mild             & 1{,}434 &  40 &  30 \\
Moderate         & 1{,}434 & 104 &  87 \\
Severe           & 1{,}434 &  22 &  17 \\
Proliferative DR & 1{,}434 &  28 &  33 \\
\midrule
Total            & 7{,}170 & 366 & 366 \\
\bottomrule
\end{tabular}
\end{table}

The following metrics are reported. \emph{Accuracy} is the overall fraction of
correctly graded images. \emph{Per-class precision, recall, and F1-score} quantify
grade-specific performance, and are aggregated as both \emph{macro} averages
(unweighted mean across classes, sensitive to minority-class performance) and
\emph{weighted} averages (support-weighted mean, dominated by frequent classes).
The \emph{quadratic weighted kappa} measures agreement with the reference grades
while penalising errors in proportion to the squared distance between the predicted
and true grade, and is therefore the most appropriate single summary for an ordinal
grading task. Discriminative ability is reported through the multiclass area under
the ROC curve (AUC), in both one-vs-rest (OVR) and one-vs-one (OVO) formulations,
together with per-class OVR-AUC.

\subsection{Training Dynamics}
\label{subsec:training_dynamics}

The two-phase training schedule is summarised in
Table~\ref{tab:training_summary}. The frozen-backbone phase reached a best
validation QWK of $0.8877$, after which fine-tuning the upper backbone further
improved the best validation QWK to $0.9059$ before early stopping was triggered.
The corresponding training curves (loss, accuracy, and validation QWK per epoch)
are shown in Figure~\ref{fig:training_curves}.

\begin{table}[H]
\centering
\caption{Two-phase training summary for the SMOTE model. Each phase restored the
weights of the epoch with the highest validation QWK.}
\label{tab:training_summary}
\begin{tabular}{lcc}
\toprule
 & Phase 1 (frozen) & Phase 2 (fine-tune) \\
\midrule
Trainable backbone layers & 0 / 238   & 109 / 238 \\
Max.\ epochs               & 60        & 50        \\
Best epoch (restored)      & 56        & 24        \\
Best validation QWK        & 0.8877    & 0.9059    \\
\bottomrule
\end{tabular}
\end{table}

% \begin{figure}[H]
%   \centering
%   \includegraphics[width=\linewidth]{figures/training_curves_frozen.png}
%   \includegraphics[width=\linewidth]{figures/training_curves_finetune.png}
%   \caption{Training and validation curves for the frozen (top) and fine-tuning
%   (bottom) phases: accuracy, loss, and validation QWK per epoch.}
%   \label{fig:training_curves}
% \end{figure}

\subsection{Classification Performance}
\label{subsec:classification_results}

\subsubsection{Overall Performance}
\label{subsubsec:overall_performance}

The headline test-set metrics are reported in Table~\ref{tab:overall_metrics}. The
model attained an overall accuracy of $0.822$ and a quadratic weighted kappa of
$0.8762$, the latter indicating substantial agreement with the reference grades on
the ordinal scale. The gap between the weighted F1-score ($0.819$) and the macro
F1-score ($0.639$) reflects the residual difficulty on the minority severity grades,
which contribute little to support-weighted averages but heavily to the unweighted
macro average.

\begin{table}[H]
\centering
\caption{Overall classification metrics on the test set ($n=366$).}
\label{tab:overall_metrics}
\begin{tabular}{lc}
\toprule
Metric & Value \\
\midrule
Accuracy                   & 0.822 \\
Quadratic Weighted Kappa   & 0.8762 \\
Macro precision            & 0.665 \\
Macro recall               & 0.630 \\
Macro F1-score             & 0.639 \\
Weighted F1-score          & 0.819 \\
\bottomrule
\end{tabular}
\end{table}

\subsubsection{Per-Class Performance}
\label{subsubsec:perclass_performance}

Table~\ref{tab:classification_report} reports the per-class classification report.
A high degree of agreement with the reference grades was achieved (QWK $= 0.8762$),
and near-perfect separation of healthy from diseased eyes was obtained, with the
No-DR class classified at a recall of $0.985$ and a one-vs-rest ROC-AUC of $0.9949$.
The principal sources of error were found to lie among adjacent severity grades: the
Mild class was most frequently confused with the Moderate class, and the Severe class
was most frequently confused with the Moderate and Proliferative DR classes. Because
the great majority of misclassifications occurred between neighbouring grades rather
than between distant ones, a high quadratic weighted kappa was retained despite a
more modest overall accuracy of $0.822$. The lowest per-class performance was
recorded for the Severe class (F1 $= 0.375$; recall $= 0.353$), which constituted
both the rarest class and the class subjected to the greatest degree of synthetic
oversampling. It was therefore concluded that, although the conservative SMOTE
strategy improved the balance of the training distribution, the discrimination of
the rarest and ordinally intermediate severity grades remained the principal
limitation of the model.

\begin{table}[H]
\centering
\caption{Per-class classification report on the test set.}
\label{tab:classification_report}
\begin{tabular}{lcccc}
\toprule
Class & Precision & Recall & F1-score & Support \\
\midrule
No DR            & 0.970 & 0.985 & 0.978 & 199 \\
Mild             & 0.486 & 0.600 & 0.537 &  30 \\
Moderate         & 0.717 & 0.759 & 0.737 &  87 \\
Severe           & 0.400 & 0.353 & 0.375 &  17 \\
Proliferative DR & 0.750 & 0.455 & 0.566 &  33 \\
\midrule
Accuracy         & --    & --    & 0.822 & 366 \\
Macro Average    & 0.665 & 0.630 & 0.639 & 366 \\
Weighted Average & 0.824 & 0.822 & 0.819 & 366 \\
\bottomrule
\end{tabular}
\end{table}

\subsubsection{Confusion Structure}
\label{subsubsec:confusion}

The full confusion matrix is reported as raw counts in
Table~\ref{tab:confusion_raw} and as row-normalised proportions (i.e.\ per-class
recall on the diagonal) in Table~\ref{tab:confusion_norm}; the same information is
visualised in Figure~\ref{fig:confusion_matrices}. The confusion structure confirms
the ordinal nature of the errors. The No-DR class was almost perfectly recovered
($196/199$ correct), with its only leakage being three images graded as Mild. The
Mild, Severe, and Proliferative classes each exchanged the majority of their errors
with their ordinal neighbours, and in particular the Moderate class acted as the
dominant ``attractor'' for misclassified minority-grade images: $10$ of $30$ Mild,
$8$ of $17$ Severe, and $8$ of $33$ Proliferative DR images were assigned to it.
Off-diagonal mass between distant grades (e.g.\ No DR vs.\ Proliferative DR) was
negligible.

\begin{table}[H]
\centering
\caption{Confusion matrix (raw counts) on the test set. Rows are true grades,
columns are predicted grades.}
\label{tab:confusion_raw}
\begin{tabular}{lccccc}
\toprule
True $\backslash$ Pred & No DR & Mild & Mod. & Sev. & Prolif. \\
\midrule
No DR            & 196 &  3 &  0 & 0 &  0 \\
Mild             &   2 & 18 & 10 & 0 &  0 \\
Moderate         &   4 & 11 & 66 & 4 &  2 \\
Severe           &   0 &  0 &  8 & 6 &  3 \\
Proliferative DR &   0 &  5 &  8 & 5 & 15 \\
\bottomrule
\end{tabular}
\end{table}

\begin{table}[H]
\centering
\caption{Row-normalised confusion matrix on the test set. Diagonal entries equal
the per-class recall. Rows may not sum to exactly $1.000$ due to rounding.}
\label{tab:confusion_norm}
\begin{tabular}{lccccc}
\toprule
True $\backslash$ Pred & No DR & Mild & Mod. & Sev. & Prolif. \\
\midrule
No DR            & 0.985 & 0.015 & 0.000 & 0.000 & 0.000 \\
Mild             & 0.067 & 0.600 & 0.333 & 0.000 & 0.000 \\
Moderate         & 0.046 & 0.126 & 0.759 & 0.046 & 0.023 \\
Severe           & 0.000 & 0.000 & 0.471 & 0.353 & 0.176 \\
Proliferative DR & 0.000 & 0.152 & 0.242 & 0.152 & 0.455 \\
\bottomrule
\end{tabular}
\end{table}

% \begin{figure}[H]
%   \centering
%   \includegraphics[width=0.49\linewidth]{figures/test_confusion_matrix_raw.png}
%   \includegraphics[width=0.49\linewidth]{figures/test_confusion_matrix_normalized.png}
%   \caption{Test-set confusion matrices: raw counts (left) and row-normalised
%   proportions (right).}
%   \label{fig:confusion_matrices}
% \end{figure}

\subsubsection{Discriminative Ability (ROC-AUC)}
\label{subsubsec:auc}

The threshold-independent discriminative ability of the model is summarised in
Table~\ref{tab:auc}, and the per-class one-vs-rest ROC curves are shown in
Figure~\ref{fig:roc_curves}. The one-vs-rest macro-averaged AUC was $0.9310$ and the
weighted AUC was $0.9611$, indicating that the softmax scores ranked images by
severity well even where the $\arg\max$ decision was incorrect. Per-class OVR-AUC
exceeded $0.91$ for every grade and reached $0.9949$ for No DR, confirming that the
limited recall on the minority grades is primarily a decision-threshold and
class-prior effect rather than a failure of the learned representation to separate
the grades.

\begin{table}[H]
\centering
\caption{Multiclass ROC-AUC on the test set. OVR: one-vs-rest; OVO: one-vs-one.}
\label{tab:auc}
\begin{tabular}{lc}
\toprule
AUC variant & Value \\
\midrule
OVR macro      & 0.9310 \\
OVR weighted   & 0.9611 \\
OVR micro      & 0.9660 \\
OVO macro      & 0.8764 \\
OVO weighted   & 0.9225 \\
\midrule
\multicolumn{2}{l}{\emph{Per-class (OVR)}} \\
\quad No DR            & 0.9949 \\
\quad Mild             & 0.8882 \\
\quad Moderate         & 0.9330 \\
\quad Severe           & 0.9200 \\
\quad Proliferative DR & 0.9186 \\
\bottomrule
\end{tabular}
\end{table}

% \begin{figure}[H]
%   \centering
%   \includegraphics[width=0.75\linewidth]{figures/test_roc_curves.png}
%   \caption{Per-class one-vs-rest ROC curves on the test set.}
%   \label{fig:roc_curves}
% \end{figure}

\subsection{Explainability Results}
\label{subsec:xai_results}

To assess whether the model's predictions were supported by clinically meaningful
image evidence, Grad-CAM~\cite{selvaraju2017gradcam} saliency maps were generated
from the final convolutional layer of the EfficientNetB0 backbone. Heatmaps were
produced for representative correctly classified and misclassified test images
across all five severity grades; a representative selection is shown in
Figure~\ref{fig:gradcam_grid}.

\subsubsection{Qualitative Saliency Analysis}
\label{subsubsec:xai_qualitative}

For correctly graded diseased images, the activated regions were generally
co-located with retinal structures relevant to DR grading, and the spatial extent
of the activated region tended to grow with predicted severity, consistent with the
broader and more numerous lesions present at higher grades. For No-DR images, the
activations were diffuse and concentrated around the major vascular arcades and the
optic disc rather than on focal lesions, in line with the absence of pathology.
Inspection of the misclassified cases was consistent with the confusion structure
reported in Section~\ref{subsubsec:confusion}: images confused between adjacent
grades typically exhibited saliency over genuine retinal regions but at a severity
boundary, rather than attending to background or artefactual regions.

% \begin{figure}[H]
%   \centering
%   \includegraphics[width=\linewidth]{figures/gradcam_examples_grid.png}
%   \caption{Grad-CAM overlays for representative test images across the five
%   severity grades. Each panel shows the input fundus image with the Grad-CAM
%   heatmap superimposed; warmer colours indicate regions that contributed more
%   strongly to the predicted grade.}
%   \label{fig:gradcam_grid}
% \end{figure}

\subsubsection{Quantitative Localisation}
\label{subsubsec:xai_quantitative}

A quantitative evaluation of saliency quality (e.g.\ the Pointing
Game~\cite{zhang2016excitation} and intersection-over-union of the thresholded
heatmap with lesion regions) requires pixel-level lesion annotations, which are not
available for the APTOS~2019 dataset: APTOS provides image-level severity labels
only. Consequently, the explainability analysis presented here is qualitative.
Quantitative localisation evaluation is left to future work and would require a
lesion-annotated dataset such as IDRiD~\cite{porwal2018idrid}, against which the
top-ranked Grad-CAM regions could be scored for overlap with annotated
microaneurysms, haemorrhages, and exudates.

\subsection{Impact of Class Imbalance Handling}
\label{subsec:imbalance_results}

%% ----------------------------------------------------------------------------
%% TODO (authors): the SMOTE + Focal-loss configuration must be run to
%% completion. Fill every "[to supply]" cell in Table~\ref{tab:imbalance_ablation}
%% and complete the prose marked below once the focal-loss test metrics,
%% confusion matrix, and AUC are available.
%% ----------------------------------------------------------------------------

This section compares the two class-imbalance strategies investigated in this work:
(i) SMOTE oversampling of the training split (the model analysed above), and
(ii) SMOTE combined with a sparse categorical focal loss
($\gamma = 2.0$, no per-class $\alpha$ weighting, so as not to double-compensate for
the already-balanced SMOTE training distribution). The focal loss down-weights
easy, already-correct examples and concentrates the gradient on hard examples, and
was motivated by the under-performance of the rarest (Severe) grade reported in
Section~\ref{subsubsec:perclass_performance}. Because both configurations share an
identical backbone, head, augmentation pipeline, training schedule, and evaluation
protocol, the comparison isolates the effect of the loss function on both
classification performance and the qualitative explanations.

Table~\ref{tab:imbalance_ablation} reports the two configurations side by side on
the identical test set. \emph{[The focal-loss column is to be completed once that
experiment is run to completion; the accompanying prose interpreting the direction
and magnitude of the change, and any effect on the minority-grade recall and on the
Grad-CAM localisation, will be added at that point.]}

\begin{table}[H]
\centering
\caption{Effect of the class-imbalance strategy on test-set performance. Both
configurations use SMOTE-balanced training data and differ only in the loss
function. Focal-loss values are pending completion of that experiment.}
\label{tab:imbalance_ablation}
\begin{tabular}{lcc}
\toprule
Metric & SMOTE & SMOTE + Focal loss \\
\midrule
Accuracy                  & 0.822  & 0.8169 \\
Macro F1-score            & 0.639  & 0.6136 \\
Weighted F1-score         & 0.819  & 0.8106 \\
Quadratic Weighted Kappa  & 0.8762 & 0.8843 \\
AUC (OVR macro)           & 0.9310 & 0.9154 \\
\midrule
\multicolumn{3}{l}{\emph{Per-class recall}} \\
\quad No DR            & 0.985 & 0.974874 \\
\quad Mild             & 0.600 & 0.366667 \\
\quad Moderate         & 0.759 & 0.839080 \\
\quad Severe           & 0.353 & 0.352941 \\
\quad Proliferative DR & 0.455 & 0.454545 \\
\bottomrule
\end{tabular}
\end{table}
